\documentclass[12pt,a4paper]{exam}
\usepackage[utf8]{inputenc}
\usepackage{amsmath,amssymb,amsthm}
\usepackage[margin=1in]{geometry}
\usepackage{graphicx}
\usepackage[colorlinks=true]{hyperref}
\pagestyle{headandfoot}
\firstpageheader{MIT OpenCourseWare}{}{\textbf{EXTREMELY HARD -- MIT LEVEL}}
\runningheader{MIT OpenCourseWare}{}{\textbf{EXTREMELY HARD -- MIT LEVEL}}
\firstpagefooter{}{Page \thepage}{}
\runningfooter{}{Page \thepage}{}

\begin{document}

\begin{center}
{\LARGE \textbf{MIT OpenCourseWare}}\\21L.004\\Major Poets
\vspace{0.5cm}
{21L.004} --- {Major Poets}
\vspace{0.3cm}
May 2026
\vspace{0.3cm}
\textbf{EXTREMELY HARD -- MIT LEVEL}
\end{center}

\begin{center}
\textbf{Time Limit: 3 hours}
\end{center}

\begin{center}
\textbf{Total Marks: 100}
\end{center}

\vspace{0.5cm}

\begin{questions}

\question[5]
Analyze the limitations of standard approaches to Prosody when applied to edge cases in 21L.004. Construct explicit counterexamples showing where intuitive reasoning fails.

\vspace{0.3cm}

\question[5]
Construct an original proof-level problem involving Poetic Form for 21L.004 (Major Poets) and provide a complete solution. Your problem should be at the level of a graduate qualifying exam.

\vspace{0.3cm}

\question[7]
Prove the fundamental theorem concerning Close Reading in the context of 21L.004 (Major Poets). State the theorem precisely, provide a complete proof, and discuss three important corollaries.

\vspace{0.3cm}

\question[3]
Construct an original proof-level problem involving Literary History for 21L.004 (Major Poets) and provide a complete solution. Your problem should be at the level of a graduate qualifying exam.

\vspace{0.3cm}

\question[2]
Construct an original proof-level problem involving Intertextuality for 21L.004 (Major Poets) and provide a complete solution. Your problem should be at the level of a graduate qualifying exam.

\vspace{0.3cm}

\question[2]
Prove the fundamental theorem concerning Prosody in the context of 21L.004 (Major Poets). State the theorem precisely, provide a complete proof, and discuss three important corollaries.

\vspace{0.3cm}

\question[5]
Analyze the limitations of standard approaches to Poetic Form when applied to edge cases in 21L.004. Construct explicit counterexamples showing where intuitive reasoning fails.

\vspace{0.3cm}

\question[3]
Analyze the limitations of standard approaches to Close Reading when applied to edge cases in 21L.004. Construct explicit counterexamples showing where intuitive reasoning fails.

\vspace{0.3cm}

\question[5]
Prove the fundamental theorem concerning Literary History in the context of 21L.004 (Major Poets). State the theorem precisely, provide a complete proof, and discuss three important corollaries.

\vspace{0.3cm}

\question[3]
Construct an original proof-level problem involving Intertextuality for 21L.004 (Major Poets) and provide a complete solution. Your problem should be at the level of a graduate qualifying exam.

\vspace{0.3cm}

\question[4]
Prove the fundamental theorem concerning Prosody in the context of 21L.004 (Major Poets). State the theorem precisely, provide a complete proof, and discuss three important corollaries.

\vspace{0.3cm}

\question[4]
Analyze the limitations of standard approaches to Poetic Form when applied to edge cases in 21L.004. Construct explicit counterexamples showing where intuitive reasoning fails.

\vspace{0.3cm}

\question[3]
Construct an original proof-level problem involving Close Reading for 21L.004 (Major Poets) and provide a complete solution. Your problem should be at the level of a graduate qualifying exam.

\vspace{0.3cm}

\question[4]
Prove the fundamental theorem concerning Literary History in the context of 21L.004 (Major Poets). State the theorem precisely, provide a complete proof, and discuss three important corollaries.

\vspace{0.3cm}

\question[2]
Construct an original proof-level problem involving Intertextuality for 21L.004 (Major Poets) and provide a complete solution. Your problem should be at the level of a graduate qualifying exam.

\vspace{0.3cm}

\question[4]
Analyze the limitations of standard approaches to Prosody when applied to edge cases in 21L.004. Construct explicit counterexamples showing where intuitive reasoning fails.

\vspace{0.3cm}

\question[3]
Construct an original proof-level problem involving Poetic Form for 21L.004 (Major Poets) and provide a complete solution. Your problem should be at the level of a graduate qualifying exam.

\vspace{0.3cm}

\question[6]
Construct an original proof-level problem involving Close Reading for 21L.004 (Major Poets) and provide a complete solution. Your problem should be at the level of a graduate qualifying exam.

\vspace{0.3cm}

\question[5]
Construct an original proof-level problem involving Literary History for 21L.004 (Major Poets) and provide a complete solution. Your problem should be at the level of a graduate qualifying exam.

\vspace{0.3cm}

\question[5]
Analyze the limitations of standard approaches to Intertextuality when applied to edge cases in 21L.004. Construct explicit counterexamples showing where intuitive reasoning fails.

\vspace{0.3cm}

\question[5]
Construct an original proof-level problem involving Prosody for 21L.004 (Major Poets) and provide a complete solution. Your problem should be at the level of a graduate qualifying exam.

\vspace{0.3cm}

\question[3]
Prove the fundamental theorem concerning Poetic Form in the context of 21L.004 (Major Poets). State the theorem precisely, provide a complete proof, and discuss three important corollaries.

\vspace{0.3cm}

\question[5]
Prove the fundamental theorem concerning Close Reading in the context of 21L.004 (Major Poets). State the theorem precisely, provide a complete proof, and discuss three important corollaries.

\vspace{0.3cm}

\question[4]
Prove the fundamental theorem concerning Literary History in the context of 21L.004 (Major Poets). State the theorem precisely, provide a complete proof, and discuss three important corollaries.

\vspace{0.3cm}

\question[3]
Prove the fundamental theorem concerning Intertextuality in the context of 21L.004 (Major Poets). State the theorem precisely, provide a complete proof, and discuss three important corollaries.

\vspace{0.3cm}

\end{questions}

\newpage
\section*{Solutions}

\textbf{Solution 1:}

The edge case analysis reveals the subtle assumptions underlying standard treatments of Prosody. By constructing explicit counterexamples, we see that the usual hypotheses cannot be weakened without loss of the theorem's conclusion.

\vspace{0.5cm}

\textbf{Solution 2:}

Solution approach: First recall the definitions and key theorems related to Poetic Form from 21L.004. Then apply the appropriate techniques to construct a rigorous argument. The solution must be self-contained and reference only the material from the course.

\vspace{0.5cm}

\textbf{Solution 3:}

The proof follows from the definitions and properties established in 21L.004. The key insight is to apply the relevant theorem and verify the hypotheses hold. Detailed solution depends on the specific formulation of the theorem.

\vspace{0.5cm}

\textbf{Solution 4:}

Solution approach: First recall the definitions and key theorems related to Literary History from 21L.004. Then apply the appropriate techniques to construct a rigorous argument. The solution must be self-contained and reference only the material from the course.

\vspace{0.5cm}

\textbf{Solution 5:}

Solution approach: First recall the definitions and key theorems related to Intertextuality from 21L.004. Then apply the appropriate techniques to construct a rigorous argument. The solution must be self-contained and reference only the material from the course.

\vspace{0.5cm}

\textbf{Solution 6:}

The proof follows from the definitions and properties established in 21L.004. The key insight is to apply the relevant theorem and verify the hypotheses hold. Detailed solution depends on the specific formulation of the theorem.

\vspace{0.5cm}

\textbf{Solution 7:}

The edge case analysis reveals the subtle assumptions underlying standard treatments of Poetic Form. By constructing explicit counterexamples, we see that the usual hypotheses cannot be weakened without loss of the theorem's conclusion.

\vspace{0.5cm}

\textbf{Solution 8:}

The edge case analysis reveals the subtle assumptions underlying standard treatments of Close Reading. By constructing explicit counterexamples, we see that the usual hypotheses cannot be weakened without loss of the theorem's conclusion.

\vspace{0.5cm}

\textbf{Solution 9:}

The proof follows from the definitions and properties established in 21L.004. The key insight is to apply the relevant theorem and verify the hypotheses hold. Detailed solution depends on the specific formulation of the theorem.

\vspace{0.5cm}

\textbf{Solution 10:}

Solution approach: First recall the definitions and key theorems related to Intertextuality from 21L.004. Then apply the appropriate techniques to construct a rigorous argument. The solution must be self-contained and reference only the material from the course.

\vspace{0.5cm}

\textbf{Solution 11:}

The proof follows from the definitions and properties established in 21L.004. The key insight is to apply the relevant theorem and verify the hypotheses hold. Detailed solution depends on the specific formulation of the theorem.

\vspace{0.5cm}

\textbf{Solution 12:}

The edge case analysis reveals the subtle assumptions underlying standard treatments of Poetic Form. By constructing explicit counterexamples, we see that the usual hypotheses cannot be weakened without loss of the theorem's conclusion.

\vspace{0.5cm}

\textbf{Solution 13:}

Solution approach: First recall the definitions and key theorems related to Close Reading from 21L.004. Then apply the appropriate techniques to construct a rigorous argument. The solution must be self-contained and reference only the material from the course.

\vspace{0.5cm}

\textbf{Solution 14:}

The proof follows from the definitions and properties established in 21L.004. The key insight is to apply the relevant theorem and verify the hypotheses hold. Detailed solution depends on the specific formulation of the theorem.

\vspace{0.5cm}

\textbf{Solution 15:}

Solution approach: First recall the definitions and key theorems related to Intertextuality from 21L.004. Then apply the appropriate techniques to construct a rigorous argument. The solution must be self-contained and reference only the material from the course.

\vspace{0.5cm}

\textbf{Solution 16:}

The edge case analysis reveals the subtle assumptions underlying standard treatments of Prosody. By constructing explicit counterexamples, we see that the usual hypotheses cannot be weakened without loss of the theorem's conclusion.

\vspace{0.5cm}

\textbf{Solution 17:}

Solution approach: First recall the definitions and key theorems related to Poetic Form from 21L.004. Then apply the appropriate techniques to construct a rigorous argument. The solution must be self-contained and reference only the material from the course.

\vspace{0.5cm}

\textbf{Solution 18:}

Solution approach: First recall the definitions and key theorems related to Close Reading from 21L.004. Then apply the appropriate techniques to construct a rigorous argument. The solution must be self-contained and reference only the material from the course.

\vspace{0.5cm}

\textbf{Solution 19:}

Solution approach: First recall the definitions and key theorems related to Literary History from 21L.004. Then apply the appropriate techniques to construct a rigorous argument. The solution must be self-contained and reference only the material from the course.

\vspace{0.5cm}

\textbf{Solution 20:}

The edge case analysis reveals the subtle assumptions underlying standard treatments of Intertextuality. By constructing explicit counterexamples, we see that the usual hypotheses cannot be weakened without loss of the theorem's conclusion.

\vspace{0.5cm}

\textbf{Solution 21:}

Solution approach: First recall the definitions and key theorems related to Prosody from 21L.004. Then apply the appropriate techniques to construct a rigorous argument. The solution must be self-contained and reference only the material from the course.

\vspace{0.5cm}

\textbf{Solution 22:}

The proof follows from the definitions and properties established in 21L.004. The key insight is to apply the relevant theorem and verify the hypotheses hold. Detailed solution depends on the specific formulation of the theorem.

\vspace{0.5cm}

\textbf{Solution 23:}

The proof follows from the definitions and properties established in 21L.004. The key insight is to apply the relevant theorem and verify the hypotheses hold. Detailed solution depends on the specific formulation of the theorem.

\vspace{0.5cm}

\textbf{Solution 24:}

The proof follows from the definitions and properties established in 21L.004. The key insight is to apply the relevant theorem and verify the hypotheses hold. Detailed solution depends on the specific formulation of the theorem.

\vspace{0.5cm}

\textbf{Solution 25:}

The proof follows from the definitions and properties established in 21L.004. The key insight is to apply the relevant theorem and verify the hypotheses hold. Detailed solution depends on the specific formulation of the theorem.

\vspace{0.5cm}

\end{document}